% ++++++++++++++++++++++++++++++++++++++++++++++++++++++++++++++++++++++++
% Autor              : Ricardo Pereira Dias <contato@ricardopdias.com.br>
% Data de criação    : 29 de Janeiro de 2019
% Referência         : https://github.com/abntex/abntex2/blob/master/tex/latex/abntex2/abntex2.cls
% ++++++++++++++++++++++++++++++++++++++++++++++++++++++++++++++++++++++++

% Opções Memoir
% 12pt       tamanho da fonte
% openright  capítulos começam em pág ímpar (insere página vazia caso preciso)
% twoside    para impressão em recto e verso. Oposto a oneside
% a4paper    % tamanho do papel.

% Opções Abntex2
% Para converter as seções pra etra maiúscula
% chapter=TITLE, section=TITLE, subsection=TITLE, subsubsection=TITLE
% Para converter as seções pra letra minúscula
% chapter=title, section=title, subsection=title, subsubsection=title
% sumario=tradicional|abnt-6027-2012
% Opções do pacote Babel
% english   idioma adicional para hifenização
% french    idioma adicional para hifenização
% spanish   idioma adicional para hifenização
% brazil    % o último idioma é o principal do documento

\documentclass[
    12pt,openright,twoside,a4paper, % Memoir
    sumario=tradicional
    brazil, % Babel
]{abntex2}

% Pacotes configurados para utilização
% ++++++++++++++++++++++++++++++++++++++++++++++++++++++++++++++++++++++++
% Autor              : Ricardo Pereira Dias <contato@ricardopdias.com.br>
% Data de criação    : 29 de Janeiro de 2019
% ++++++++++++++++++++++++++++++++++++++++++++++++++++++++++++++++++++++++

% =========================================================================
% PACOTES
% =========================================================================


\usepackage[utf8]{inputenc}  % Codificacao do documento (conversão automática dos acentos)
\usepackage{indentfirst}     % Indenta o primeiro parágrafo de cada seção.
\usepackage{color}           % Controle das cores
\usepackage{graphicx}        % Inclusão de gráficos
\usepackage{microtype}       % para melhorias de justificação
\usepackage{xcolor}
\usepackage{fancyhdr}
\usepackage{longtable}
\usepackage{colortbl}
\usepackage{lipsum}       % para geração de dummy text
\usepackage[subentrycounter,seeautonumberlist,nonumberlist=true]{glossaries}
\usepackage[brazilian,hyperpageref]{backref}     % Paginas com as citações na bibl
\usepackage[alf]{abntex2cite}    % Citações padrão ABNT


% Informações pré-definidas
%=========================================================================
% Autor              : Ricardo Pereira Dias <rpdesignerfly@gmail.com>
%-------------------------------------------------------------------------
% Data de criação    : 29 de Janeiro de 2019
% Cores: https://www.overleaf.com/learn/latex/Using_colours_in_LaTeX
%=========================================================================

%=========================================================================
% CAPA e FOLHA DE ROSTO
%=========================================================================

% Projeto de Pesquisa
% TCC
% Artigo
\newcommand{\schoolInstitution}{%
    Universidade de Raccoon City - URC
    \par
    Faculdade de Arquitetura da Informação
    \par
    Programa de Pós-Graduação
}
\newcommand{\projectType}{Trabalho Acadêmico (TCC)}
\newcommand{\projectName}{Modelo Canônico de \projectType{} com Speed Latex}
\newcommand{\projectCity}{São Carlos - SP}
\newcommand{\projectDate}{2019}
\newcommand{\projectVersion}{, v-1.0.0}
% O preambulo deve conter o tipo do trabalho, o objetivo,
% o nome da instituição e a área de concentração
\newcommand{\projectPreamble}{Modelo canônico de  \projectType{} baseado no Abntex2 em conformidade com as normas ABNT, apresentado à comunidade de usuários \LaTeX.}
\newcommand{\projectAuthor}{Ricardo Pereira Dias}
\newcommand{\projectContact}{contato@ricardopdias.com.br}

% TCC
\newcommand{\projectOrientingLabel}{Prof.}%rotulo
\newcommand{\projectOrienting}{Albert Wesker}
\newcommand{\projectCoorientingLabel}{Prof.}%rotulo
\newcommand{\projectCoorienting}{Barry Burton}

% Artigo
\newcommand{\projectForeignName}{Canonical Academic Work template in SpeedLatex} % Nome extrangeiro

%=========================================================================
% CABEÇALHO
%=========================================================================
% Aparência
\definecolor{header_rule}{RGB}{ 0, 0, 0 }
\definecolor{header_text}{RGB}{ 0, 0, 0 }
\newcommand{\headerRuleHeight}{0pt}
% Páginas ímpares
\newcommand{\headerOddLeft}{}
\newcommand{\headerOddCenter}{}
\newcommand{\headerOddRight}{\thepage}
% Páginas Pares
\newcommand{\headerEvenLeft}{\thepage}
\newcommand{\headerEvenCenter}{}
\newcommand{\headerEvenRight}{}


%=========================================================================
% RODAPÉ
%=========================================================================
% Aparência
\definecolor{footer_rule}{RGB}{ 0, 0, 0 }
\definecolor{footer_text}{RGB}{ 0, 0, 0 }
\newcommand{\footerRuleHeight}{0pt}
% Páginas ímpares
\newcommand{\footerOddLeft}{}
\newcommand{\footerOddCenter}{}
\newcommand{\footerOddRight}{}
% Páginas Pares
\newcommand{\footerEvenLeft}{}
\newcommand{\footerEvenCenter}{}
\newcommand{\footerEvenRight}{}

%=========================================================================
% CORES GERAIS
%=========================================================================
% http://erikasarti.com/html/tabela-cores/
\definecolor{black}{RGB}{0,0,0}
% \definecolor{blue}{RGB}{41,5,195}
\definecolor{blue}{RGB}{0,0,0}
\definecolor{dim_gray}{RGB}{105,105,105}
\definecolor{gray}{RGB}{128,128,128}
\definecolor{dark_gray}{RGB}{169,169,169}
\definecolor{silver}{RGB}{192,192,192}
\definecolor{light_grey}{RGB}{211,211,211}
\definecolor{gainsboro}{RGB}{220,220,220}
\definecolor{white}{RGB}{255,255,255}

\definecolor{success_text}{RGB}{ 43, 180, 54 }
\definecolor{info_text}{RGB}{ 2, 69, 156 }
\definecolor{warning_text}{RGB}{ 255, 158, 0}
\definecolor{danger_text}{RGB}{ 255, 60, 0 }

%=========================================================================
% CAPÍTULOS E SEÇÕES
%=========================================================================

% Vários exemplos de capitulos
% http://tug.ctan.org/info/MemoirChapStyles/
% https://texblog.org/2012/07/03/fancy-latex-chapter-styles/

% Para formatar todos os niveis de seção
\newcommand{\sectioningTopStyle}{\bfseries}
\newcommand{\sectioningSecStyle}{\bfseries}
\newcommand{\sectioningFinalDot}{} % Para poder adicionar um character no final da numeração


% Partes
% Declaração: Parte 1
\newcommand{\partAssertionFont}{\sffamily}
\newcommand{\partAssertionStyle}{\sectioningTopStyle}
\newcommand{\partAssertionSize}{\large}
% Título da parte
\newcommand{\partFont}{\sffamily}
\newcommand{\partStyle}{\sectioningTopStyle}
\newcommand{\partSize}{\Huge}
\definecolor{part_count}{RGB}{ 0, 0, 0 }
\definecolor{part_text}{RGB}{ 0, 0, 0 }

% Capítulos
\newcommand{\chapterFont}{\sffamily}
\newcommand{\chapterStyle}{\sectioningTopStyle}
\newcommand{\chapterSize}{\Huge}
\definecolor{chapter_count}{RGB}{ 0, 0, 0 }
\definecolor{chapter_text}{RGB}{ 0, 0, 0 }

% Seções
\newcommand{\sectionFont}{\sffamily}
\newcommand{\sectionStyle}{\sectioningSecStyle}
\newcommand{\sectionSize}{\Large}
\definecolor{section_count}{RGB}{ 0, 0, 0 }
\definecolor{section_text}{RGB}{ 0, 0, 0 }

% Sub Seções
\newcommand{\subsectionFont}{\sffamily}
\newcommand{\subsectionStyle}{\sectioningSecStyle}
\newcommand{\subsectionSize}{\large}
\definecolor{subsection_count}{RGB}{ 0, 0, 0 }
\definecolor{subsection_text}{RGB}{ 0, 0, 0 }

% Sub Sub Seções
\newcommand{\subsubsectionFont}{\sffamily}
\newcommand{\subsubsectionStyle}{\sectioningSecStyle}
\newcommand{\subsubsectionSize}{\normalsize}
\definecolor{subsubsection_count}{RGB}{ 0, 0, 0 }
\definecolor{subsubsection_text}{RGB}{ 0, 0, 0 }

% Sub Sub Sub Seções
\newcommand{\subsubsubsectionFont}{\sffamily}
\newcommand{\subsubsubsectionStyle}{\sectioningSecStyle}
\newcommand{\subsubsubsectionSize}{\normalsize}
\definecolor{subsubsubsection_count}{RGB}{ 0, 0, 0 }
\definecolor{subsubsubsection_text}{RGB}{ 0, 0, 0 }

\definecolor{paragraph_text}{RGB}{ 0, 0, 0 } % parágrafo de leis ou documentos
\definecolor{subparagraph_text}{RGB}{ 0, 0, 0 } % subparágrafo de leis ou documentos

% Corpo de texto
\definecolor{body_text}{RGB}{ 0, 0, 0 }
\definecolor{link_text}{RGB}{ 0, 0, 0 }


% Configurações do documento
\include{libraries/document-config-common}

% Fontes que acompanham o texlive
% % =========================================================================
% Autor              : Ricardo Pereira Dias <rpdesignerfly@gmail.com>
% Data de criação    : 30 de Março de 2019
% =========================================================================
% Para listar, no terminal: "fc-list : family"
% -------------------------------------------------------------------------
% Links interssantes:
% http://www.tug.dk/FontCatalogue/
% http://www.tug.dk/FontCatalogue/allfonts.html
% http://www.tug.dk/FontCatalogue/cabin/
% http://mirrors.ibiblio.org/CTAN/fonts/
% -------------------------------------------------------------------------

% code: pag

\usepackage{avant}
\usepackage[T1]{fontenc}

% % =========================================================================
% Autor              : Ricardo Pereira Dias <rpdesignerfly@gmail.com>
% Data de criação    : 30 de Março de 2019
% =========================================================================
% Para listar todas as fontes disponiveis no sistema,
% digite no terminal:
% fc-list : family
% -------------------------------------------------------------------------
% Links interssantes:
% http://www.tug.dk/FontCatalogue/
% https://en.wikibooks.org/wiki/LaTeX/Fonts
% https://pt.overleaf.com/learn/latex/Font_sizes,_families,_and_styles#Font_families
% -------------------------------------------------------------------------

% https://ctan.org/tex-archive/fonts/urw/arial

\usepackage{uarial}
\usepackage[T1]{fontenc}

% % =========================================================================
% Autor              : Ricardo Pereira Dias <rpdesignerfly@gmail.com>
% Data de criação    : 30 de Março de 2019
% =========================================================================
% OpenSans
% -------------------------------------------------------------------------
% Para listar todas as fontes disponiveis no sistema,
% fc-list : family
% -------------------------------------------------------------------------
% Links interssantes:
% http://www.tug.dk/FontCatalogue/
% http://www.tug.dk/FontCatalogue/allfonts.html
% http://www.tug.dk/FontCatalogue/cabin/
% -------------------------------------------------------------------------

% https://ctan.org/tex-archive/fonts/bera

% Bera Serif: fve
% Bera Sans: fvs
% Bera Sans Mono: fvm

\usepackage{bera}
\usepackage[T1]{fontenc}

% \include{libraries/fonts/cabin}
% % =========================================================================
% Autor              : Ricardo Pereira Dias <rpdesignerfly@gmail.com>
% Data de criação    : 30 de Março de 2019
% =========================================================================
% Para listar todas as fontes disponiveis no sistema,
% digite no terminal:
% fc-list : family
% -------------------------------------------------------------------------
% Links interssantes:
% http://www.tug.dk/FontCatalogue/
% http://www.tug.dk/FontCatalogue/firasansbook/
% -------------------------------------------------------------------------

% https://ctan.org/tex-archive/fonts/fira

\usepackage[book]{FiraSans}
\usepackage[T1]{fontenc}
\renewcommand*\oldstylenums[1]{{\firaoldstyle #1}}

% % =========================================================================
% Autor              : Ricardo Pereira Dias <rpdesignerfly@gmail.com>
% Data de criação    : 29 de Janeiro de 2019
% =========================================================================
% COMPUTER MODERN
% -------------------------------------------------------------------------
% Para listar todas as fontes disponiveis no sistema,
% digite no terminal:
% fc-list : family
% -------------------------------------------------------------------------
% Links interssantes:
% http://www.tug.dk/FontCatalogue/
% https://en.wikibooks.org/wiki/LaTeX/Fonts
% https://pt.overleaf.com/learn/latex/Font_sizes,_families,_and_styles#Font_families
% -------------------------------------------------------------------------

\usepackage{lmodern}
\usepackage[T1]{fontenc}

\include{libraries/fonts/roboto}

% Inclui as entradas do glossario
\include{templates/common-document/glossary}

\begin{document}

    % Define o idioma do documento (conforme pacotes do babel)
    \selectlanguage{brazil}

    % Retira espaço extra obsoleto entre as frases.
    \frenchspacing

    %---
    %
    % Se desejar escrever o artigo em duas colunas, descomente a linha abaixo
    % e a linha com o texto ``FIM DE ARTIGO EM DUAS COLUNAS''.
    % \twocolumn[    		% INICIO DE ARTIGO EM DUAS COLUNAS
    %
    %---

    %===============================================
    % ELEMENTOS PRÉ-TEXTUAIS
    %===============================================
    %% ----------------------------------------------
% Capa
% ----------------------------------------------

\begin{capa}%

    \center
    
    \includegraphics[scale=.5]{assets/logo.eps}

    \vfill

    \chapterFont\bfseries\LARGE\projectName

    \vfill

    \chapterFont\large\projectAuthor

    \vfill

    \large\projectCity
    \large\projectDate

    \vspace*{1cm}
\end{capa}
%
    \input{templates/common-document/pre-textuais/02-sumario}%

    %===============================================
    % ELEMENTOS TEXTUAIS
    %===============================================
    \textual

    % 
% ----------------------------------------------------------
% Introdução (exemplo de capítulo sem numeração, mas presente no Sumário)
% ----------------------------------------------------------
\chapter{Introdução}
% ----------------------------------------------------------

Este documento e seu código-fonte são exemplos de referência de uso do Speed Latex,
que faz uso da excelente classe \textsf{memoir} e da customização oferecida pelo projeto \textsf{\abnTeX}. Optou-se por utilizar a classe \textsf{\abnTeX} sem alterá-la, seguindo a ``recomendação dos autores do projeto''\footnote{\url{https://github.com/abntex/abntex2/wiki/ComoCustomizar}}.

\textbf{Atenção:} Embora faça-se uso do \textsf{\abnTeX}, o modelo "Livro Simples"\ não segue as regras da ABNT.

\section{Modelos disponíveis}

O Speed Latex disponibiliza vários modelos prontos para criação de documentos baseados em Latex.
Confira abaixo a lista de modelos disponíveis:

\begin{itemize}

    \item \textbf{abnt-academic-work} - Para Trabalhos Acadêmicos/TCC (ABNT);
    \item \textbf{abnt-article} - Para Artigos Científicos (ABNT);
    \item \textbf{abnt-research-project} - Para Projetos de Pesquisa (ABNT);
    \item \textbf{common-letter} - Para criação de Cartas;
    \item \textbf{common-document} - Para criação de Documentos;
    \item \textbf{common-book} - Para criação de Livros

\end{itemize}

\section{Organização}

Além de possibilitar a criação de projetos pré-organizados baseados em Latex, o Speed Latex
permite uma maior facilidade na personalização do documento. Isso é feito através de um arquivo de configuração existente para cada projeto, onde possível personalizar várias informações, como:

\begin{itemize}

    \item O conteúdo dos cabeçalhos e rodapés;
    \item As cores das seções e dos textos em geral;
    \item Definir informações padrões para uso no documento

\end{itemize}

\section{Atualizações}

O Speed Latex \footnote{\url{https://github.com/ricardopedias/speed-latex}} está em constante atualização. Sinta-se à vontade para conferir os releases \footnote{\url{https://github.com/ricardopedias/speed-latex/releases}} e acompanhar as atualizações.

Sinceramente espera-se que cada atualização do Speed Latex aprimore a qualidade do trabalho que você produzirá, de modo que o principal esforço seja concentrado no conteúdo produzido e não na formatação de documentos.
%
    % % ---
% Capitulo que faz uso de elementos do glossario
% ---
\section{Orientações a respeito de glossários}

Você pode definir as entradas do glossário no início do texto. Recomenda-se o
uso de um arquivo separado a ser inserido ainda no preâmbulo. Veja orientações
sobre inclusão de arquivos na \autoref{sec-include}.

No decorrer do texto, use os termos do glossário como na frase:

\begin{citacao}
Esta frase usa a palavra \gls{componente} e o plural de \glspl{filho}, ambas
definidas no glossário como filhas da entrada \gls{pai}. \Gls{equilibrio}
exemplifica o uso de um termo no início da frase. O software \gls{abntex2} é
escrito em \gls{latex}, que é definido no glossário como
\emph{`\glsdesc*{latex}'}.
\end{citacao}


A frase acima foi produzida com:

\begin{verbatim}
Esta frase usa a palavra \gls{componente} e o plural de \glspl{filho}, ambas
definidas no glossário como filhas da entrada \gls{pai}. \Gls{equilibrio}
exemplifica o uso de um termo no início da frase. O software \gls{abntex2} é
escrito em \gls{latex}, que é definido no glossário como
\emph{`\glsdesc*{latex}'}.
\end{verbatim}

Opcionalmente, incorpore todas as palavras do glossário de uma única vez ao
documento com o comando:

\begin{verbatim}
   \glsaddall
\end{verbatim}

A impressão efetiva do glossário é dada com:

\begin{verbatim}
   \printglossaries
\end{verbatim}

A impressão do glossário incorpora o número das páginas em que as entradas foram
citadas. Isso pode ser removido adicionando-se a opção \texttt{nonumberlist} em:

\begin{verbatim}
\usepackage[nonumberlist,style=index]{glossaries}%
\end{verbatim}

% ---
\section{Compilar um documento com glossário}
\label{sec-compilar-glossario}
% ---

Para compilar um documento \LaTeX\ com glossário use:

\begin{verbatim}
   pdflatex ARQUIVO_PRINCIPAL.tex
   bibtex ARQUIVO_PRINCIPAL.aux
   makeindex ARQUIVO_PRINCIPAL.idx
   makeindex ARQUIVO_PRINCIPAL.nlo -s nomencl.ist -o ARQUIVO_PRINCIPAL.nls
   makeglossaries ARQUIVO_PRINCIPAL.aux
   pdflatex ARQUIVO_PRINCIPAL.tex
   pdflatex ARQUIVO_PRINCIPAL.tex
\end{verbatim}

O comando \texttt{makeglossaries} é um aplicativo Perl instalado
automaticamente pelas distribuições MacTeX, TeX Live e MiKTeX. Geralmente
usuários de Linux e de Mac OS X já possuem o interpretador Perl\footnote{O Perl
é uma linguagem de programação de scripts muito utilizada pela comunidade de
software livre. Veja o site do projeto em \url{http://www.perl.org/}.} instalado
e configurado e nenhuma configuração adicional é necessária.

Usuários de Windows, por outro lado, precisam instalar a ferramenta Perl para
que seja possível usar \texttt{makeglossaries}. Por sorte isso é simples. Para
obter a instalação do Perl para seu sistema operacional visite \url{http://www.perl.org/get.html}.

Alternativamente ao aplicativo Perl \texttt{makeglossaries}, é possível usar o
aplicativo \texttt{makeglossariesgui}\footnote{O título do aplicativo no CTAN
é \textit{Java GUI alternative to makeglossarires script}.}, que possui uma
interface gráfica baseada em Java. Para isso, consulte
\url{http://www.ctan.org/pkg/makeglossariesgui}. Funciona em Windows,
Linux e Mac OS X.

% ---
\section{Configuração de glossários}
% ---

O pacote \textsf{glossaries}, usado na produção dos glossários deste exemplo,
possui diversas configurações. É possível alterar o estilo da impressão do
glossário, criar campos adicionais, usar diversos glossários em
arquivos separados. Para isso e outras informações, consulte a documentação do
pacote \textsf{glossaries}: \url{http://www.ctan.org/pkg/glossaries}.

Consulte também o livro da WikiBooks sobre a produção de glossários:
\url{http://en.wikibooks.org/wiki/LaTeX/Glossary}.


\subsection{Estilos do glossário}

O pacote \textsf{glossaries} traz dezenas de estilos pré-definidos de
glossários. Eles estão disponíveis no capítulo 15 do manual do pacote
\cite{talbot2012}. O capítulo 16 contém instruções sobre como criar um estilo
personalizado.

Os estilos podem ser alterados com:

\begin{verbatim}
   \setglossarystyle{altlisthypergroup}
\end{verbatim}

O estilo \texttt{index} é ideal para construção de glossários com diversos
níveis hierárquicos do tipo pai-filho. Já o modelo \texttt{altlisthypergroup} é
mais adequado para glossários sem hierarquias. Teste também o modelo
\texttt{tree}.

Se desejar um único estilo de glossário padrão no documento, alternativamente
inclua a opção \texttt{style} nas opções da classe, do
seguinte modo:

\begin{verbatim}
   \usepackage[style=index]{glossaries}
\end{verbatim}

% ---
\section{Problemas com a ordem das palavras?}
% ---

Este exemplo do \abnTeX\ utiliza a ferramenta \texttt{makeindex} -- padrão das
distribuições \LaTeX\ mais comuns -- para ordenar as entradas do glossário.
Porém, essa ferramenta não possui opções de \textit{collation} e não funciona
bem para palavras escritas em idiomas que não sejam inglês.
Por isso, pode acontecer que letras acentuadas e outros caracteres
internacionais sejam ordenados de forma incorreta, como no exemplo (palavras não
necessariamente presentes na língua portuguesa):

\begin{alineas}
 \item Amor: ...
 \item Aviar: ...
 \item Avião: ...
 \item Aço: ...
\end{alineas}

Por sorte, é possível substituir o uso do \texttt{makeindex}
pelo \texttt{xindy}\footnote{\url{http://www.xindy.org/}}. Para isso, faça o
seguinte:

\begin{alineas}
  \item Certifique-se de que o Xindy esteja instalado. Em um terminal, digite:
  \texttt{xindy --version}\footnote{Caso o Xindy não esteja presente no sistema, é necessário
    instalá-lo. Usuários Linux Debian/Ubuntu podem usar: \texttt{sudo
    apt-get install xindy}. Usuários Windows e Mac podem acessar a página do
    Xindy, baixá-lo e instalá-lo.};
  \item No código LaTeX, ainda no preâmbulo, inclua a seguinte opção ao pacote glossaries:
  \begin{verbatim}
  \usepackage[xindy={language=portuguese},nonumberlist=true]{glossaries}
  \end{verbatim}
  \item Compile o glossário normalmente, conforme a
  \autoref{sec-compilar-glossario}.
\end{alineas}

    % \input{templates/common-document/textuais/02-secao-comandos}%
    % \input{templates/common-document/textuais/03-secao-cabecalhos-rodapes}%
    % 


\section{Mais exemplos no Modelo Canônico de Trabalhos Acadêmicos}

Este modelo de artigo é limitado em número de exemplos de comandos, pois são
apresentados exclusivamente comandos diretamente relacionados com a produção de
artigos.

Para exemplos adicionais de \abnTeX\ e \LaTeX, como inclusão de figuras,
fórmulas matemáticas, citações, e outros, consulte o documento
\citeonline{abntex2modelo}.
%
    % 


\section{Consulte o manual da classe \textsf{abntex2}}

Consulte o manual da classe \textsf{abntex2} \cite{abntex2classe} para uma
referência completa das macros e ambientes disponíveis.
%

    % ---
    % Finaliza a parte no bookmark do PDF, para que se inicie o bookmark na raiz
    % ---
    \bookmarksetup{startatroot}%

    % 
    %===================================================================================================
    % ELEMENTOS TEXTUAIS
    %===================================================================================================

    % \pagestyle{meuestilo}
    % \textual

    % ------------------------------------------------------------------------
    % Conclusão
    % ------------------------------------------------------------------------
    \chapter*[Considerações finais]{Considerações finais}
    \addcontentsline{toc}{chapter}{Considerações finais}

    \lipsum[31-33]
%


    %===============================================
    % ELEMENTOS PÓS-TEXTUAIS
    %===============================================
    \postextual

    % ----------------------------------------------------------
    % Referências bibliográficas
    % ----------------------------------------------------------
    \bibliography{templates/common-document/pos-textuais/01-references}
    % \input{templates/common-document/pos-textuais/02-glossario}
    % \input{templates/common-document/pos-textuais/03-apendices}
    % \cftinserthook{toc}{AAA}
    % 
% ----------------------------------------------
% Anexos
% ----------------------------------------------
\begin{anexosenv}

    % ---
    \chapter{Cras non urna sed feugiat cum sociis natoque penatibus et magnis dis
    parturient montes nascetur ridiculus mus}
    % ---

    \lipsum[31]

    \vspace*{1.5cm}

\end{anexosenv}

    % \input{templates/common-document/pos-textuais/05-agradecimentos}

\end{document}
